\chapter{The Logical Topology and Compactness}
\label{sec:topology}
% ===================================================================

\section{The Logical Topology}

\begin{definition}[Logical Topology]
\label{def:log-top}
Let $\GSLT$ be a GSLT with context-decorated HML $\HML(\ctx)$.
For each formula $\phi \in \HML(\ctx)$, define the
\emph{extension} of $\phi$:
\[
  \llb\phi\rrb \;=\; \{ [P] \in \terms(\GSLT)/{\bisim} \mid P \sat \phi \}.
\]
The \emph{logical topology} $\mathcal{O}_{\HML}$ on
$\terms(\GSLT)/{\bisim}$ is the topology generated by the
extensions $\{ \llb\phi\rrb \mid \phi \in \HML(\ctx) \}$ as a
subbasis.
\end{definition}

\begin{remark}[Stone, not Scott]
\label{rmk:stone}
Earlier drafts called this the \emph{Scott topology} and the resulting
space a spectral space --- sober, $T_0$, quasi-compact.
That is the right description for a subbasis of positive observations
and it is the wrong description for this one.
$\HML(\ctx)$ has negation.
Definition~\ref{def:cdHML} carries it, the namespace logic of \chSci\ uses
it, and Conjecture~\ref{thm:compact-consensus} quantifies over
$\neg\phi$ explicitly.
So $\llb \neg\phi \rrb$ is the complement of $\llb \phi \rrb$, every
subbasic open is also closed, and the space is zero-dimensional.
By the Hennessy--Milner property distinct bisimulation classes are
separated by some formula, hence by a clopen set, so the space is
Hausdorff and its points are closed.

The setting is therefore \emph{Stone} duality and not domain theory.
The specialisation order, which carries the content in a Scott
topology, is discrete here and carries none; ``sober and $T_0$''
understates the separation badly.
This is not a question of which citation to attach.
The two settings answer differently when asked what compactness of a
population means.
In the Stone setting $X_\pop$ compact says that
$\terms(\GSLT)/{\bisim}$ restricted to the population is the Stone
space of the Boolean algebra of $\HML(\ctx)$-definable sets --- a
profinite limit of finite quotients, each one the population as seen
at a bounded modal depth.

That is the same structure \S\ref{sci:sec:lattice} of \chSci\ puts on
the hypothesis space, where the ultrametric comes from agreement up to
depth and the balls are the levels of the same inverse system.
The two parts have been using one object under two descriptions.
The logical metric $d_{\HML}$ of Part~\ref{part:machinery} is the
ultrametric that induces this topology, and the metric balls are
exactly the basic clopens.
\end{remark}

\section{Population Subspaces}

\begin{definition}[Latent Subspace]
Given a population $\pop \subseteq \terms(\GSLT)/{\bisim}$, the
\emph{latent topological subspace} of $\pop$ is the subspace
\[
  X_\pop \;=\; \overline{\pop} \;\subseteq\; \terms(\GSLT)/{\bisim}
\]
with the subspace topology inherited from $\mathcal{O}_{\HML}$.
The closure is taken in the logical topology.
\end{definition}

\begin{remark}
The latent subspace is ``latent'' in the sense that it is not
constructed deliberately by the agents.
It is the topological shadow cast by the population's collective
behavior --- the smallest closed set in the logical topology that
contains all the agents.
An agent inside $\pop$ cannot directly observe $X_\pop$;
it can only probe it indirectly through experiments.
\end{remark}

\section{Compactness as Consensus}

\begin{definition}[Compact Population]
A population $\pop$ is \emph{compact} if its latent subspace
$X_\pop$ is compact in the logical topology $\mathcal{O}_{\HML}$:
every open cover of $X_\pop$ by extensions of HML formulae has a
finite subcover.
\end{definition}

\begin{conjecture}[Compactness implies consensus]
\label{thm:compact-consensus}
Let $\pop$ be a compact population in an interactive GSLT $\GSLT$.
Then $\pop$ supports consensus: there is no infinite oscillation
on any Boolean question expressible in $\HML(\ctx)$.

More precisely: for any formula $\phi \in \HML(\ctx)$, if the
population is divided into two communities indefinitely alternating
between $\phi$ and $\neg\phi$, then the alternation must terminate
in finite time.
\end{conjecture}

This was stated as a theorem in earlier drafts and the argument given
for it does not work.
The argument, and both of the obvious ways to save it, are set out
below, because the failure is more instructive than the statement and
because a reader who is going to rely on this ought to see exactly what
is missing.

\begin{remark}[The argument that was given, and why it fails]
\label{rmk:consensus-attempt}
The proof sketch defined a family of formulae of increasing modal
depth,
\[
  \phi_1 = \langle 0 \rangle\top, \quad
  \phi_2 = \langle 1 \rangle\langle 0 \rangle\top, \quad
  \phi_3 = \langle 0 \rangle\langle 1 \rangle\langle 0 \rangle\top,
  \quad \ldots
\]
with $0$ and $1$ labelling the two communities' states, so that
$\phi_n$ holds of an agent that has completed \emph{at least} $n$
oscillation steps.
It then claimed that $\{ \llb \phi_n \rrb \}_{n \in \NN}$ is an open
cover of $X_\pop$ with no finite subcover.

It is not.
Satisfying $\phi_{n+1}$ entails satisfying $\phi_n$, so the family is
\emph{decreasing}:
$\llb \phi_1 \rrb \supseteq \llb \phi_2 \rrb \supseteq \cdots$.
Hence $\bigcup_{n \le N} \llb \phi_n \rrb = \llb \phi_1 \rrb$ for
every $N$, and an agent that has oscillated more than $N$ times lies
\emph{inside} that union rather than outside it.
There is no failure of finite subcovering and therefore no
contradiction with compactness.

The natural repair is to reindex, taking $\psi_n = \neg\phi_n$ ---
``has completed fewer than $n$ steps'' --- so that the family is
increasing.
That family is an open cover of $X_\pop$ precisely when every agent
oscillates only finitely often, which is the conclusion.
The repair assumes what it sets out to prove.

The other natural repair runs compactness in its finite-intersection
form.
The $\llb \phi_n \rrb$ are closed as well as open, by
Remark~\ref{rmk:stone}, and the family has the finite intersection
property whenever some agent has oscillated $N$ times for each $N$.
Compactness then yields
$\bigcap_n \llb \phi_n \rrb \neq \emptyset$ --- an agent satisfying
every $\phi_n$, which is to say an infinite oscillator.
This is a correct argument and it establishes the opposite of the
conjecture.
Compactness is a condition guaranteeing that limits are attained, and a
perpetual disagreement is exactly such a limit.
\end{remark}

\begin{remark}[What the conjecture is probably missing]
\label{rmk:consensus-missing}
The third argument having gone the wrong way is the informative one.
Compactness in the logical topology is a \emph{richness} condition: by
Stone duality it says that every finitely satisfiable set of formulae
is realised by some point of $X_\pop$, so the population's closure
omits no consistent type.
Nothing about richness makes a process stop.

Termination is not a topological property of the population; it is a
budgetary property of the process that revises theories.
The framework has the missing hypothesis and it is not in this chapter.
Every revision is an assay and every assay is paid for
(\S\ref{sci:sec:search} of \chSci), so if each alternation between
$\phi$ and $\neg\phi$ costs a bounded-below quantity from a
bounded-above endowment, the alternation terminates, and it does so for
reasons that have nothing to do with the topology.
Compactness would then contribute the second half of the statement
rather than the first: it bounds the depth at which the surviving
disagreement can be located, and hence the convergence \emph{time},
once termination is available from the budget.

Two routes out, then.
Add the cost hypothesis and prove the conjecture as a statement about
priced revision, in which case it belongs in
Part~\ref{part:ecology} and not here.
Or keep it topological and weaken the conclusion from ``the
alternation terminates'' to ``the alternation is confined to a
subspace of bounded logical depth'', which may be provable as stated
and is weaker than what the rest of this part uses.
The first looks more likely to be true and more useful.
Until one of them is done this is a conjecture, and the two results
downstream of it --- coherence clusters as maximal compact
subpopulations, and the reading of non-compactness as permanent
disagreement --- should be read as depending on it.
\end{remark}

\begin{remark}[If it holds, the proof delivers the protocol]
Conjecture~\ref{thm:compact-consensus} is an existence claim:
it says consensus \emph{must} be reached, without specifying how.
A proof of compactness for a \emph{specific} population subspace is
constructive, however: it identifies the finite subcover, which encodes
the maximum depth of oscillation and therefore the convergence time.
The proof of compactness would be the consensus protocol.
This is the sense in which the topology was intended to do
computational work, and it is worth stating because it is the payoff
that makes the conjecture worth settling.
\end{remark}

\begin{remark}[Non-compact populations]
On the intended reading, a non-compact population may contain
irresolvable oscillations --- communities that never agree ---
and non-compactness is the topological signature of
\emph{fundamental disagreement}: not a disagreement that more
communication or better arguments could resolve, but one that is
structurally permanent.
In the context of the massive population, non-compactness at a given
scale would mean the cluster at that scale cannot function as a
coherent agent.
Remark~\ref{rmk:consensus-missing} is the reason this is stated in the
conditional.
Under Stone duality non-compactness says the population omits a
consistent type, which is a statement about \emph{gaps} in the
population rather than about deadlock within it, and the two readings
have not been reconciled.
\end{remark}

% ===================================================================
